\chapter{Wahlanzeige}
\quellseite{11}

In Folge der Aufforderung im „Lutheraner“ vom 22. December 1903, Candidaten für die vacante Professur am Gymnasium in Milwaukee aufzustellen, sind folgende Personen nominirt worden:

1. P. Dr. H. Wente, Halfway, Mich.

2. P. Phil. Spengling, Saginaw, Wis.

3. P. W. Dallmann, New York.

4. Prof. M. Walter, Hawthorne, N. Y.

5. P. Leo Schmidtke, Chippewa Falls, Wis.

6. P. Theo. Steege, Lyons, Iowa.

7. P. Friedr. Meyer, Jefferson City, Mo.

8. P. Herm. Spechard, Saginaw, Mich.

9. P. S. Gänßle, Corning, Mo.

10. P. Herm. Eggers, Watertown, Wis.

11. P. J. R. Gräbner, Salt Lake City, Utah.

12. P. G. Birkmann, Seboy, Tex.

13. P. Alb. Bartling, Waterford, Wis.

14. P. Dr. H. W. F. Wolläger, Hartford, Conn.

15. P. F. S. Bünger, St. Louis, Mo.

16. P. G. Schlegel, Dorsen, Ill.

17. P. Rob. Rißling, Madison, Wis.

18. P. Phil. Wilhelm, Manville, Wis.

19. Prof. H. Jürgensen, St. Paul, Minn.

20. Dir. H. W. F. Albrecht, Milwaukee, Wis.

21. P. Theo. Schlüter, Courtland, Minn.

Diese Candidaten stehen nun während der nächsten vier Wochen vom Datum dieser Anzeige ab, das ist, bis zum 1. März, vor der Synode, damit Gelegenheit gegeben werde, etwaige Proteste einzureichen. Sobald die Zeit verflossen ist, wird die Wahl stattfinden.

\begin{unterschrift}
Fort Wayne, Ind., den 2. Februar 1904.\\
C. Groß sen.,\\
Secretär des Wahlcollegiums
\end{unterschrift}
